\scanpage{35}
\mspage{26}
% archive sheet number: 35
% inv: text | top | Ainsi une figure indexée
Ainsi, une figure indexée \inserted{par $I$} définit une application
% inv: display | top | varphi I to M
\[
\varphi:I\longrightarrow M
\]
% inv: text | upper third | qui est croissante
qui est \uline{croissante} pour $\leq$ sur $M$, \deleted{Quelles} et $F$ est déterminée par $\varphi$ \deleted{\unsure{comme}} avec
% inv: display | upper third | tilde F equals varphi(I)
\[
\widetilde F=\varphi(I).
\]
% inv: text | upper middle | Quelles conditions
Quelles conditions sur $\varphi$, pour que ça corresponde bien à une figure indexée ? Il faut

% inv: text | upper middle | a) varphi injectif
 a) $\varphi$ injectif, et \deleted{$=$} \inserted{induit} un isom. d'ens. ordonnés $I\simeq\varphi(I)$

% inv: text | middle | et b) varphi(I) est une figure
et\qquad b) $\varphi(I)$ est une figure. \deleted{\unsure{Je dis que}} \deleted{Cela signifie} \deleted{deux choses}

% inv: display | middle | 1 degrees forall i j
\[
1^\circ)\quad\forall i,j\in I,\quad X_i\unsure{\gtrless|\bowtie}X_j\quad\deleted{\unsure{\text{i.e.}}}\quad\text{et}
\]
% inv: text | middle | 2 degrees varphi(I)
$2^\circ)$\quad $\varphi(I)$ est une partie fermée de $M$.

% inv: text | middle lower | La première condition
La première condition revient à celle-ci ; \deleted{qui implique aussi \unsure{a)}}
% inv: display | middle lower | b) restriction
\[
b)\quad\forall i\in I,\quad\varphi|I_{\leq i}:I_{\leq i}\xrightarrow{\sim}\widetilde X_i
\]
% inv: display | middle lower | c) Si i not equal j
\[
c)\quad\text{Si }i\ne j,\quad X_i\mathrel{|\circ|}X_j
\]

% inv: text | lower third | En d'autres termes
En d'autres termes, \unsure{je} dis $X\longrightarrow X_i$,

% inv: text | lower third | Proposition
\uline{Proposition}\quad Pour que $\varphi:I\longrightarrow M$, définisse une figure indexée par $I$, il f et s que

% inv: text | lower third | a) varphi croissante
 a) $\varphi$ \deleted{injectif} croissante (ds $M,\leq$) et \uline{injective}.

% inv: text | lower third | b) application induite
 b) $\forall i\in I$, l'application induite
% inv: display | bottom third | tilde i equals I le i
\[
\deleted{\unsure{\mathrm{de}}}\quad\widetilde i=I_{\leq i}\xrightarrow{\sim}\widetilde X_i=M_{\leq X_i}
\]
% inv: text | bottom third | soit bijective
soit \deleted{\unsure{bijective} (\unsure{compatible} \unsure{toujours|tjrs} \unsure{aux} \unsure{relations} \unsure{d'ordre}} $=$ un isom. d'ensembles ordonnés)

% inv: text | bottom | c) Si i not equal j
 c) Si $i\ne j$ (i.e $i\mathrel{|\circ|}j$ pour la structure de magasin trivial sur $I$ !) \unsure{:} $X_i\mathrel{|\circ|}X_j$.
